% The study's paper: one LaTeX document, yours to write. Load any package with \usepackage, and include images
% from the study's folder with \includegraphics (a path from paper/, e.g. figures/chart.png, or from the study).
% Bind every number to the record: \lab{result:<id>#metrics.<metric>.value}, \labtab{result:<id>},
% \labfig{result:<id>}, \labclaim{claim:<id>}, \labcite{<source>}{<quote>}. \labsection{<name>} puts in the text of a
% section openlab writes from the record (data, methods, literature, results, limitations, reproducibility,
% checklist). A number typed in the prose is a finding: the draft shows it, the signed build refuses it.
\documentclass[11pt]{article}
\usepackage[margin=2.4cm]{geometry}
\usepackage{graphicx,booktabs}

\title{<<title>>}
\author{<<authors>>}
\date{}

\begin{document}
\maketitle

\begin{abstract}
What was asked, what was done, and what was found, in a few sentences.
\end{abstract}

\section{Introduction}
The question and why it matters.

\section{Data}
Where the data come from and what they hold. % \labsection{data} puts openlab's account of the data here

\section{Methods}
The model, the evaluation and the baselines. % \labsection{methods} puts openlab's account of the plan here

\section{Results}
What the runs found. % \labtab{result:<id>} tables a result; \labsection{results} lists the current ones

\section{Discussion}
What the results mean, and what they do not show.

\section{Limitations}
% \labsection{limitations} lists what the record knows: the trials, the holdout, the assumptions at risk

\section{Reproducibility}
\labsection{reproducibility}

\end{document}
